\documentclass[11pt]{amsart}
\usepackage{amsmath,amssymb,amsthm}
\usepackage[margin=1.15in]{geometry}
\usepackage[hidelinks]{hyperref}

\newtheorem{theorem}{Theorem}[section]
\newtheorem{lemma}[theorem]{Lemma}
\newtheorem{corollary}[theorem]{Corollary}
\newtheorem{proposition}[theorem]{Proposition}
\theoremstyle{definition}
\newtheorem{definition}[theorem]{Definition}
\theoremstyle{remark}
\newtheorem{remark}[theorem]{Remark}

\newcommand{\arxiv}[1]{\href{https://arxiv.org/abs/#1}{arXiv:\nolinkurl{#1}}}
\newcommand{\N}{\mathcal{N}}
\newcommand{\Q}{Q}
\newcommand{\F}{\mathbb{F}}
\newcommand{\Z}{\mathbb{Z}}
\newcommand{\Hred}{\widetilde{H}}
\newcommand{\cl}[2]{N_{#1}[#2]}
\newcommand{\C}{\operatorname{C}}
\setlength{\emergencystretch}{2.5em}

\begin{document}

\title[Collapsibility and homology of \v{C}ech complexes of cubes]{Collapsibility numbers and high-degree homology of
\v{C}ech complexes of hypercube graphs}
\author{Yitzchak Shmalo}
\date{26 November 2026}
\address{Einstein Institute of Mathematics, The Hebrew University of Jerusalem, Jerusalem, Israel}
\email{yspennstate@gmail.com}
\subjclass[2020]{55U10, 05E45, 05C12, 55N31}
\keywords{\v{C}ech complex, hypercube graph, collapsibility number, $d$-collapsible, set-pair system, independent
dominating set}

\begin{abstract}
Let $\N(\Q_n,r)$ be the \v{C}ech complex of the hypercube graph $\Q_n$ at integer scale $r$. We show that its
collapsibility number is $2^{k+1}-1$ at every even scale $r=2k$ and every $n\ge k+1$, the value proposed by Adams,
Shukla and Singh. At odd scale $r=2m-3$ with $m\ge3$ we construct, for each $h$ with $2\le h\le\lfloor\log_2(m+1)\rfloor$,
a free direct summand of rank at least $2^h-1$ in $\Hred_{q}(\N(\Q_n,r);\Z)$, $q=2^m-2^{m-h}-2$, for every $n\ge m$. The
summand for $h=3$ lies in a degree outside the set proposed by Adams, Shukla and Singh for every odd $r\ge11$, and it
forces the collapsibility number above the proposed odd-scale value $3\cdot2^k-1$. Together with an upper bound of
$2^{k+2}-1$, the collapsibility number at scale $2k+1$ divided by $2^{k+2}$ tends to $1$, uniformly in $n\ge k+2$. For $k=1$ and $m\ge7$ the construction proves the first part of a conjecture of Conant and
Thistlethwaite: the complex $\Theta(\mathrm{Cube}(m,1))$ is not homotopy equivalent to a wedge of spheres of a single dimension.
\end{abstract}

\maketitle

\section{Introduction}

For a connected graph $G$ and a real $r\ge0$, the \v{C}ech complex $\N(G,r)$ is the nerve of the closed balls of radius
$r/2$ about the vertices of $G$ in its geodesic realization. Adams, Shukla and Singh \cite{ASS} began the study of these
complexes for the hypercube graph $\Q_n$ (written $\mathbb{I}_n$ there), the graph on $\F_2^n$ in which two words are
adjacent when they differ in one coordinate. They determined the homotopy types at scales $r\le 2$, proved results at
$r=3$, and recorded homology computed by machine in a table. Section~7 of \cite{ASS} lists open questions. Two of them
concern us. Question~7.2 asks: for $r=2k$ even and $n\ge k+1$, is $\Hred_i(\N(\Q_n,r))\ne0$ if and only if
$i\in\{r,\,2(2^k-1)\}$; and for $r=2k+1$ odd and $n\ge k+2$, is $\Hred_i(\N(\Q_n,r))\ne0$ if and only if
$i\in\{r,\,3\cdot2^k-2\}$? Question~7.4 asks whether the collapsibility number of $\N(\Q_n,2k)$ is $2^{k+1}-1$ for
$n\ge k+1$, and whether that of $\N(\Q_n,2k+1)$ is $3\cdot2^k-1$ for $n\ge k+2$. Adams, Shukla and Singh prove the first
value at $k=1$ \cite[Lemma~7.3]{ASS} and the second at $k=1$ for $n\ge4$ \cite[Theorem~5.9]{ASS}.

Throughout, $\Hred_*$ is reduced simplicial homology. Write $\C(X)$ for the collapsibility number of a simplicial
complex $X$, the least $d$ such that $X$ is $d$-collapsible
in the sense of Wegner \cite{W}.

\begin{theorem}\label{thm:even}
For every $k\ge0$ and $n\ge k+1$, $\C(\N(\Q_n,2k))=2^{k+1}-1$.
\end{theorem}

\begin{theorem}\label{thm:rank}
Let $m\ge3$, $N=2^m$ and $2\le h\le\lfloor\log_2(m+1)\rfloor$. For every $n\ge m$, the group
$\Hred_{q}(\N(\Q_n,2m-3);\Z)$ with $q=N-N/2^h-2$ contains a free direct summand of rank at least $2^h-1$.
\end{theorem}

\begin{theorem}\label{thm:odd}
For every $k\ge1$ and $n\ge k+2$, with $m=k+2$, $N=2^m$ and $h=\lfloor\log_2(m+1)\rfloor$,
\[
N-\frac{N}{2^h}-1\;\le\;\C(\N(\Q_n,2k+1))\;\le\;N-1 .
\]
In particular $\C(\N(\Q_n,2k+1))/2^{k+2}\to1$ as $k\to\infty$, uniformly in $n\ge k+2$.
\end{theorem}

For $h=2$ the degree in Theorem~\ref{thm:rank} is $3\cdot2^k-2$ with $k=m-2$, the second degree in Question~7.2, and the
lower bound of Theorem~\ref{thm:odd} with $h=2$ is the proposed value $3\cdot2^k-1$. For $m\ge7$ the choice $h=3$ is
available, and the following is immediate.

\begin{corollary}\label{cor:questions}
Let $k\ge5$ and $n\ge k+2$. Then $\Hred_{7\cdot2^{k-1}-2}(\N(\Q_n,2k+1);\Z)\ne0$, a degree different from $2k+1$ and from
$3\cdot2^k-2$; and $\C(\N(\Q_n,2k+1))\ge7\cdot2^{k-1}-1>3\cdot2^k-1$. So at every odd scale $r\ge11$ the answer to the
``only if'' part of the odd case of Question~7.2 is negative, and so is the answer to the odd case of Question~7.4.
\end{corollary}

The complex $\N(\Q_m,2m-3)$ is the complex $\Theta(\mathrm{Cube}(m,1))$ of Conant and Thistlethwaite \cite{CT}, as
Adams, Shukla and Singh observe \cite{ASS}. Conant and Thistlethwaite computed its rational homology for $m\le5$; at
$m=5$ they found homology in degrees $22$ and $24$ \cite[Theorem~5]{CT}. With Lemma~\ref{lem:retract} below, which holds
with any coefficients, the class in degree $24$ gives $\Hred_{24}(\N(\Q_n,7);\mathbb{Q})\ne0$ and $\C(\N(\Q_n,7))\ge25>3\cdot2^3-1$ for every $n\ge5$, so at $r=7$ the
odd cases of both questions are already answered by their computation. Corollary~\ref{cor:questions} answers them at
every odd $r\ge11$. The cases $r=5$ and $r=9$ are not settled by these arguments.

Conant and Thistlethwaite conjecture that $\Theta(\mathrm{Cube}(n,k))$ is not homotopy equivalent to a wedge of spheres
of a single dimension for $0<k<n-2$ and $n$ large, and that its nonzero homology spans a range of degrees that grows with $n$
\cite[Conjecture~2]{CT}. For $k=1$ and $m\ge7$, Theorem~\ref{thm:rank} with $h=2$ and $h=3$ gives nonzero integral
homology of $\N(\Q_m,2m-3)=\Theta(\mathrm{Cube}(m,1))$ in the two degrees $3\cdot2^{m-2}-2$ and $7\cdot2^{m-3}-2$, which
proves the first part of the conjecture for $k=1$ and $m\ge7$ (the case $m=5$ is their own computation). The degrees
$2^m-2^{m-h}-2$, $2\le h\le\lfloor\log_2(m+1)\rfloor$, span a range of length at least $2^{m-2}-2^{m-\lfloor\log_2(m+1)\rfloor}$,
which grows with $m$.

Two related results concern the same complexes. The inclusion $\N(\Q_n,r)\hookrightarrow\N(\Q_n,r+1)$ is
null-homotopic for every $n\ge1$ and $r\ge0$ \cite{Sh1}, which answers Question~7.5 of \cite{ASS}; it is a statement about maps between scales and
says nothing about the homology of a single complex. For each coordinate $p$-face $F$ of $\Q_n$, the reduced homology of
$\N(\Q_p,2p-3)$ embeds as a split summand of that of $\N(\Q_n,2p-3)$, independently over the faces \cite{FC}; that
result transports classes from $\Q_p$ to $\Q_n$ but proves nonvanishing on $\Q_p$ only for $p=3$. Theorem~\ref{thm:rank}
supplies the nonvanishing for every $p\ge3$, and with \cite{FC} it gives a free direct summand of rank at least
$(2^h-1)\,2^{n-m}\binom{n}{m}$ in $\Hred_q(\N(\Q_n,2m-3);\Z)$. Lemma~\ref{lem:retract}, the result of \cite{FC} and
Theorem~\ref{thm:rank} are homology propagation
results for \v{C}ech complexes in the sense of Adams and Virk, who ask for such results in \cite[Question~8.8]{AV}.

The upper bounds come from a criterion of Lew \cite{Lew} together with the skew Hamming set-pair bound of Gao and Zhao
\cite{GZ}; Selvaraja \cite[Section~5]{Sel} uses the same two results for Vietoris--Rips complexes of hypercube graphs,
and our argument for the upper bounds follows that model. The lower bounds come from spheres built on independent
dominating sets of $\Q_m$ (the fibres of a syndrome map) and from retractions onto them.

\section{Preliminaries}

\subsection{The complexes}
For a vertex $v$ of $\Q_n$ write $\cl{j}{v}$ for the set of words at Hamming distance at most $j$ from $v$, and
$\bar v$ for the complement of $v$. By \cite[Lemma~2.2]{ASS}, $\N(\Q_n,r)$ is the simplicial complex on $\F_2^n$ whose
faces are the subsets of the sets
\[
\cl{k}{v}\ (v\in\F_2^n)\quad\text{if } r=2k,\qquad
\cl{k}{u}\cup\cl{k}{w}\ (uw\text{ an edge})\quad\text{if } r=2k+1 .
\]
We call these sets the generating sets. We use only this description.

\begin{lemma}\label{lem:retract}
Let $1\le m\le n$ and $r\ge0$. The coordinate inclusion $\iota\colon\F_2^m\to\F_2^n$, $x\mapsto(x,0)$, and the projection
$\pi\colon\F_2^n\to\F_2^m$ onto the first $m$ coordinates are simplicial maps $\N(\Q_m,r)\to\N(\Q_n,r)$ and
$\N(\Q_n,r)\to\N(\Q_m,r)$, and $\pi\iota$ is the identity. Hence $\Hred_*(\N(\Q_m,r))$ is a direct summand of
$\Hred_*(\N(\Q_n,r))$, with any coefficients.
\end{lemma}

For $r=3$ this is \cite[Lemma~5.7]{ASS}, and the proof is the same.

\begin{proof}
Both maps do not increase Hamming distance. Hence $\iota(\cl{j}{v})\subseteq\cl{j}{\iota v}$ and
$\pi(\cl{j}{v})\subseteq\cl{j}{\pi v}$. At even scale this is all we need. At odd scale, $\iota$ carries an edge to an
edge. For an edge $uw$ of $\Q_n$, either $\pi u\pi w$ is an edge of $\Q_m$, and then
$\pi(\cl{k}{u}\cup\cl{k}{w})\subseteq\cl{k}{\pi u}\cup\cl{k}{\pi w}$, or $\pi u=\pi w$, and then the image lies in
$\cl{k}{\pi u}$, which lies in $\cl{k}{\pi u}\cup\cl{k}{z}$ for any neighbour $z$ of $\pi u$. Finally $\pi\iota=\mathrm{id}$.
\end{proof}

\subsection{Collapsibility}
A complex is $d$-collapsible if it can be reduced to the void complex by removing, one at a time, a face $\sigma$ of
dimension at most $d-1$ that lies in a unique maximal face, together with all faces containing $\sigma$ \cite{W}. A
$d$-collapsible complex $X$ has $\Hred_i(X)=0$ for $i\ge d$ \cite{W}; so
\begin{equation}\label{eq:homlower}
\Hred_q(X;\Z)\ne0\ \Longrightarrow\ \C(X)\ge q+1 .
\end{equation}

For the upper bounds we use Lew's criterion \cite[Theorem~4]{Lew}. For a complex $X$, let $d'(X)$ be the largest $t$
for which there are maximal faces $\sigma_1,\dots,\sigma_{t+1}$ of $X$ and vertices $v_1,\dots,v_t$ with
$v_i\notin\sigma_i$ for $i\le t$ and $v_i\in\sigma_j$ for $i<j\le t+1$. Then $X$ is $d'(X)$-collapsible, so
$\C(X)\le d'(X)$.

\subsection{The set-pair bound}
\begin{theorem}[Gao--Zhao {\cite[Theorem~2.2]{GZ}}]\label{thm:GZ}
Let $\tau\ge0$ and $n\ge\tau+1$, and let $a_1,b_1,\dots,a_M,b_M\in\F_2^n$ satisfy $d(a_i,b_i)\ge\tau+1$ for every $i$ and
$d(a_i,b_j)\le\tau$ whenever $i<j$. Then $M\le2^{\tau+1}$.
\end{theorem}

The theorem in \cite{GZ} is stated over any alphabet; we use the binary case. It is used only for the upper bounds in
Theorems~\ref{thm:even} and~\ref{thm:odd}; the lower bounds and Theorem~\ref{thm:rank} do not depend on it.

\section{Even scales}

\begin{proof}[Proof of Theorem~\ref{thm:even}]
\emph{Lower bound.} In $\Q_{k+1}$ the set $\cl{k}{v}$ is the whole vertex set except $\bar v$. So the faces of
$\N(\Q_{k+1},2k)$ are the proper subsets of $\F_2^{k+1}$, and the complex is the boundary of a simplex on $2^{k+1}$
vertices, a sphere of dimension $2^{k+1}-2$; this is noted in \cite{ASS} in the discussion of their Table~1. By Lemma~\ref{lem:retract}, $\Hred_{2^{k+1}-2}(\N(\Q_n,2k);\Z)\ne0$ for
every $n\ge k+1$, and \eqref{eq:homlower} gives $\C\ge2^{k+1}-1$.

\emph{Upper bound.} Let $X=\N(\Q_n,2k)$ with $n\ge k+1$. Every maximal face of $X$ is a generating set $\cl{k}{c}$, and
$\cl{k}{c}$ misses $\bar c$, so $X$ is not a simplex. Take $\sigma_1,\dots,\sigma_{t+1}$ and $v_1,\dots,v_t$ as in the
definition of $d'(X)$, with $\sigma_j=\cl{k}{c_j}$, and choose $v_{t+1}\notin\sigma_{t+1}$. Then $d(v_i,c_i)\ge k+1$
for $i\le t+1$ and $d(v_i,c_j)\le k$ for $i<j\le t+1$. Theorem~\ref{thm:GZ} with threshold $k$ gives $t+1\le2^{k+1}$.
Hence $\C(X)\le d'(X)\le2^{k+1}-1$.
\end{proof}

For $k=1$ this is \cite[Lemma~7.3]{ASS}, and for $k=0$ the complex is a discrete set of at least two points.

\section{Spheres in $\N(\Q_m,2m-3)$}

Throughout this section $m\ge3$, $N=2^m$, $V=\F_2^m$ and $X_m=\N(\Q_m,2m-3)$.

\begin{lemma}\label{lem:faces}
A set $\sigma\subseteq V$ is a face of $X_m$ if and only if $V\setminus\sigma$ contains an edge of $\Q_m$.
\end{lemma}

\begin{proof}
The scale is $2k+1$ with $k=m-2$. Let $uw$ be an edge. A vertex $z$ lies outside $\cl{m-2}{u}\cup\cl{m-2}{w}$ exactly
when $d(z,u)\ge m-1$ and $d(z,w)\ge m-1$. These two distances differ by one, so they are $m-1$ and $m$ in some order,
and $z\in\{\bar u,\bar w\}$. So the generating sets are the complements of the edges $\bar u\bar w$, and every edge
arises this way.
\end{proof}

\begin{lemma}\label{lem:sphere}
Let $S\subseteq V$ be independent and dominating in $\Q_m$, and let $W=V\setminus S$. The full subcomplex of $X_m$ on
$W$ is the boundary $\Sigma_W$ of the simplex on $W$, a sphere of dimension $|W|-2$.
\end{lemma}

\begin{proof}
The complement of $W$ is $S$, which contains no edge, so $W$ is not a face. For $w\in W$ some $s\in S$ is adjacent to
$w$, so the complement of $W\setminus\{w\}$ contains the edge $sw$, and $W\setminus\{w\}$ is a face.
\end{proof}

\begin{lemma}\label{lem:retraction}
Let $S$ and $W$ be as in Lemma~\ref{lem:sphere}. Suppose $S'\subseteq W$ and $\varphi\colon S\to S'$ is a bijection with
$s\varphi(s)$ an edge for every $s$, such that the subgraph of $\Q_m$ induced on $S\cup S'$ is the perfect matching
$\{s\varphi(s)\}$. Then $f\colon V\to W$, $f(s)=\varphi(s)$ for $s\in S$ and $f(w)=w$ for $w\in W$, is a simplicial
map $X_m\to\Sigma_W$, and it restricts to the identity on $\Sigma_W$.
\end{lemma}

\begin{proof}
By Lemma~\ref{lem:sphere} it is enough to show that $f(\sigma)\ne W$ for every face $\sigma$. Suppose $f(\sigma)=W$.
A vertex of $W\setminus S'$ has only itself as a preimage, so $W\setminus S'\subseteq\sigma$, and for each $s\in S$ the
face $\sigma$ contains $s$ or $\varphi(s)$. Thus $V\setminus\sigma\subseteq S\cup S'$ and it contains at most one
endpoint of each matching edge. Since the induced subgraph on $S\cup S'$ is that matching, $V\setminus\sigma$ is
independent, and by Lemma~\ref{lem:faces} $\sigma$ is not a face.
\end{proof}

\subsection*{Syndrome fibres}
Fix $h$ with $2\le h\le\lfloor\log_2(m+1)\rfloor$, so $2^h-1\le m$. Let $H$ be the $h\times m$ matrix over $\F_2$ whose
first $2^h-1$ columns are the nonzero vectors of $\F_2^h$, each once, and whose remaining $m-2^h+1$ columns all equal
one fixed nonzero vector $c_0$. Write $H_j$ for the $j$th column. Call a column \emph{simple} if no other column equals
it. The simple columns are all nonzero vectors other than $c_0$ when $m>2^h-1$, and all of them when $m=2^h-1$; as
$h\ge2$, they span $\F_2^h$. For $s\in\F_2^h$ put $S_s=\{x\in V:Hx=s\}$ and $W_s=V\setminus S_s$. Since $H$ has rank
$h$, $|S_s|=N/2^h$ for every $s$.

\begin{lemma}\label{lem:fibres}
\begin{enumerate}
\item[(a)] Each $S_s$ is independent and dominating, so $\Sigma_s:=\Sigma_{W_s}$ is a sphere in $X_m$ of dimension
$q=N-N/2^h-2$.
\item[(b)] If $H_j$ is simple and $t=H_j$, the subgraph induced on $S_s\cup S_{s+t}$ is the perfect matching
$x\leftrightarrow x+e_j$, $x\in S_s$.
\end{enumerate}
\end{lemma}

\begin{proof}
(a) Flipping coordinate $i$ changes $Hx$ by $H_i\ne0$, so no two words of $S_s$ are adjacent. If $Hx\ne s$, then
$s-Hx=H_i$ for some $i$, and $x+e_i\in S_s$. (b) Words in one fibre are not adjacent. If $x\in S_s$ and $x+e_i\in
S_{s+t}$, then $H_i=t$, so $i=j$ because $H_j$ is simple.
\end{proof}

Write $z_s\in\Hred_q(X_m;\Z)$ for the image of a fundamental class of $\Sigma_s$. If $H_j$ is simple and $t=H_j$,
Lemmas~\ref{lem:retraction} and \ref{lem:fibres}(b), with $S=S_s$, $S'=S_{s+t}$ and $\varphi(x)=x+e_j$, give a simplicial
retraction $f_{s,j}\colon X_m\to\Sigma_s$.

\begin{lemma}\label{lem:matrix}
Let $H_j$ be simple, $t=H_j$, and $u\in\F_2^h$. Identify $\Hred_q(\Sigma_s;\Z)$ with $\Z$ by the fundamental class. Then $(f_{s,j})_*$ carries the class of
$\Sigma_u$ to $1$ if $u=s$, to $\pm1$ if $u=s+t$, and to $0$ otherwise.
\end{lemma}

\begin{proof}
On $\Sigma_s$ the map is the identity. On $W_{s+t}$ it is a bijection onto $W_s$: it exchanges $S_s\subseteq W_{s+t}$
with $S_{s+t}\subseteq W_s$ and fixes the rest. So it carries $\Sigma_{s+t}$ isomorphically onto $\Sigma_s$, of degree
$\pm1$. For $u\notin\{s,s+t\}$ we have $S_s\cup S_{s+t}\subseteq W_u$, and $f_{s,j}(W_u)=W_u\setminus S_s$, a proper
subset of $W_s$ since it misses $S_u\ne\emptyset$. So $f_{s,j}$ maps $\Sigma_u$ into a simplex of $\Sigma_s$, and
$(f_{s,j})_*z_u=0$.
\end{proof}

\begin{proof}[Proof of Theorem~\ref{thm:rank}]
The Cayley graph of $\F_2^h$ with respect to the simple columns is connected, so it has a spanning tree rooted at $0$.
For $s\ne0$ let $p(s)=s+t_s$ be the parent of $s$, with $t_s=H_{j(s)}$ simple, and let $f_s=f_{s,j(s)}$. List the
$2^h-1$ nonzero $s$ in order of nondecreasing depth in the tree. Let $\Psi\colon\Z^{2^h-1}\to\Hred_q(X_m)$ send the basis vector
$e_u$ to $z_u$, and let $\Phi\colon\Hred_q(X_m)\to\Z^{2^h-1}$ have components $(f_s)_*$, after identifying
$\Hred_q(\Sigma_s)$ with $\Z$. By Lemma~\ref{lem:matrix}, the matrix of $\Phi\Psi$ has entry $1$ at $(s,s)$, entry $\pm1$
at $(s,p(s))$ when $p(s)\ne0$, and zeros elsewhere. Parents come earlier in the order, so this matrix is unitriangular
and invertible over $\Z$. Hence $\Psi$ is a split injection, and $\Hred_q(X_m;\Z)$ has a free summand of rank $2^h-1$.
Lemma~\ref{lem:retract} carries it to $\N(\Q_n,2m-3)$ for every $n\ge m$.
\end{proof}

\section{Odd scales}

\begin{proof}[Proof of Theorem~\ref{thm:odd}]
\emph{Lower bound.} Here $2k+1=2m-3$, and $n\ge k+2=m$. Theorem~\ref{thm:rank} with $h=\lfloor\log_2(m+1)\rfloor$ gives
$\Hred_{N-N/2^h-2}(\N(\Q_n,2k+1);\Z)\ne0$, and \eqref{eq:homlower} gives the lower bound.

\emph{Upper bound.} Let $X=\N(\Q_n,2k+1)$ with $n\ge k+2$. A generating set $\cl{k}{u}\cup\cl{k}{w}$ misses $\bar u$,
which is at distance $n\ge k+2$ from $u$ and $n-1\ge k+1$ from $w$; so $X$ is not a simplex. Take
$\sigma_j=\cl{k}{u_j}\cup\cl{k}{w_j}$ ($j\le t+1$) and $v_1,\dots,v_t$ as in the definition of $d'(X)$, and choose
$v_{t+1}\notin\sigma_{t+1}$. Since $v_j\notin\sigma_j$, both $d(v_j,u_j)$ and $d(v_j,w_j)$ are at least $k+1$; they
differ by one, so one of them is at least $k+2$. Call that endpoint $c_j$. For $i<j$, $v_i\in\sigma_j$ lies within $k$
of $u_j$ or of $w_j$, hence within $k+1$ of both, and $d(v_i,c_j)\le k+1$. Theorem~\ref{thm:GZ} with threshold $k+1$
(note $n\ge k+2$) gives $t+1\le2^{k+2}$, so $\C(X)\le2^{k+2}-1=N-1$.

\emph{Limit.} $h=\lfloor\log_2(m+1)\rfloor$ satisfies $2^{h}>(m+1)/2$, so the two bounds differ by at most
$N\cdot 2/(m+1)$, and the ratio to $N=2^{k+2}$ tends to $1$.
\end{proof}

\begin{proof}[Proof of Corollary~\ref{cor:questions}]
For $k\ge5$, $m=k+2\ge7$ and $h=3$ is allowed. The degree is $2^{k+2}-2^{k-1}-2=7\cdot2^{k-1}-2$, which differs from
$3\cdot2^k-2=6\cdot2^{k-1}-2$ and exceeds $2k+1$. The collapsibility bound follows from \eqref{eq:homlower}.
\end{proof}

\begin{remark}
At $k=1$ the bounds of Theorem~\ref{thm:odd} are $5$ and $7$; for $n\ge4$ the value is $5$ by
\cite[Theorem~5.9]{ASS}, and at $n=3$ the complex $\N(\Q_3,3)$ has reduced homology only in degree $4$, of rank $3$
over $\F_p$ by the script cited below. For
$2\le k\le4$ only $h=2$ is available and the lower bound is $3\cdot2^k-1$. At $k=3$ this is not sharp, by the bound
$\C(\N(\Q_n,7))\ge25$ from \cite{CT} recalled in the introduction. The rank bound of Theorem~\ref{thm:rank} is not
sharp: for $m=4$ and $m=5$ the rank in degree $q$ is $7$ and $10$ (over $\F_p$, $p=1000003$, by the script cited below;
for $m=5$ also \cite[Theorem~5]{CT}), against $2^h-1=3$. The exact odd-scale values remain open.
\end{remark}

\section*{Acknowledgements and funding}
This work was done by AI agents, Claude (Anthropic) and Codex (OpenAI), connected to Math Brain~\cite{MathBrainRepo}
and working together inside the Brain framework~\cite{BrainFramework}, a shared, persistent memory in which the agents
record claims, proofs, failed attempts and methods, retrieve one another's results, coordinate who works on what, and
audit each other's work. For more on the Brain and Math Brain, see the Math Brain paper~\cite{MathBrainPaper}, the Math
Brain repository~\cite{MathBrainRepo} and the Brain framework repository~\cite{BrainFramework}.

The research presented in this paper was supported by the European Research Council (ERC) under the European Union's
Horizon Europe research and innovation programme (grant agreement No.~101041711), by the Simons Foundation as part of
the Collaboration on the Mathematical and Scientific Foundations of Deep Learning, by Heights Labs, by Convex Nexus
Capital, by the Israel Science Foundation (grant number 2258/19), and by the Israel Science Foundation (ISF Grant
4101/25).

\section*{Declaration of competing interest}
The author holds equity in Heights Labs and in Convex Nexus Capital and has a management role at Convex Nexus Capital.
These entities are included in the funding declaration above.

\section*{Data availability}
No data were generated or analysed in this work. Python scripts that check the syndrome-fibre construction of
Section~4 for $6\le m\le9$ and compute the reduced Betti numbers of $\N(\Q_m,2m-3)$ for $m\le5$ are at
\url{https://github.com/yspennstate/cech-hypercube-odd-scales/tree/v1.4}.

\section*{Declaration of generative AI and AI-assisted technologies}
This paper is part of a research project on AI proving problems. Codex agents (OpenAI) produced the proofs of
Theorems~\ref{thm:even}, \ref{thm:rank} and \ref{thm:odd}, including the syndrome-fibre construction, and Claude agents
(Anthropic) checked the proofs and the references and wrote the manuscript. The
questions are due to Adams, Shukla and Singh~\cite{ASS}. The upper bounds follow Selvaraja's use~\cite{Sel} of Lew's
criterion~\cite{Lew} with the set-pair bound of Gao and Zhao~\cite{GZ}. The case $r=7$ rests on the computation of
Conant and Thistlethwaite~\cite{CT}. The author directed the work and is responsible
for the content of the paper and for decisions concerning its dissemination.

\begin{thebibliography}{99}
\bibitem{ASS} H.~Adams, S.~Shukla and A.~Singh, \emph{\v{C}ech complexes of hypercube graphs}, Homology Homotopy Appl.
\textbf{27} (2025), no.~1, 83--105, \href{https://doi.org/10.4310/HHA.2025.v27.n1.a6}{doi:10.4310/HHA.2025.v27.n1.a6};
\arxiv{2212.05871}. Numbered results are cited as in the arXiv version.
\bibitem{AV} H.~Adams and \v{Z}.~Virk, \emph{Lower bounds on the homology of Vietoris--Rips complexes of hypercube
graphs}, Bull. Malays. Math. Sci. Soc. \textbf{47} (2024), no.~3, Paper No.~72,
\href{https://doi.org/10.1007/s40840-024-01663-x}{doi:10.1007/s40840-024-01663-x}; \arxiv{2309.06222}. Numbered results
are cited as in the arXiv version.
\bibitem{CT} J.~Conant and O.~Thistlethwaite, \emph{Boolean formulae, hypergraphs and combinatorial topology},
Topology Appl. \textbf{157} (2010), no.~16, 2449--2461,
\href{https://doi.org/10.1016/j.topol.2010.08.016}{doi:10.1016/j.topol.2010.08.016}; \arxiv{0808.0739}.
\bibitem{FC} Y.~Shmalo, \emph{Independent face classes in \v{C}ech complexes of hypercube graphs}, 2026,
\href{https://doi.org/10.5281/zenodo.24556515}{doi:10.5281/zenodo.24556515}.
\bibitem{GZ} G.~Gao and R.~Zhao, \emph{Tight bound for the skew Hamming set-pair problem}, preprint (2026),
\arxiv{2607.17261}.
\bibitem{Lew} A.~Lew, \emph{Collapsibility of simplicial complexes of hypergraphs}, Electron. J. Combin. \textbf{26}
(2019), no.~4, Paper No.~4.10; \arxiv{1810.11802}. Numbered results are cited as in the journal version.
\bibitem{Sh1} Y.~Shmalo, \emph{A one-step nullhomotopy for \v{C}ech complexes of hypercube graphs}, 2026,
\href{https://doi.org/10.5281/zenodo.24548858}{doi:10.5281/zenodo.24548858}.
\bibitem{Sel} S.~Selvaraja, \emph{Collapsibility of Vietoris--Rips complexes of hypercube graphs in codimension two},
preprint (2026), \arxiv{2611.20820}.
\bibitem{W} G.~Wegner, \emph{$d$-collapsing and nerves of families of convex sets}, Arch. Math. (Basel) \textbf{26}
(1975), 317--321.
\bibitem{MathBrainRepo} Y.~Shmalo, \emph{Math Brain: public snapshot, code and tests}, 2026,
\url{https://github.com/yspennstate/math-brain-public/tree/5427734deee9be643d6dfbdd5e7d4baeeefe4a4a}.
\bibitem{MathBrainPaper} Y.~Shmalo, \emph{Math Brain: shared, auditable memory for multi-agent mathematical
research}, 2026, \href{https://doi.org/10.5281/zenodo.24150837}{doi:10.5281/zenodo.24150837},
\url{https://github.com/yspennstate/math-brain-public/tree/5427734deee9be643d6dfbdd5e7d4baeeefe4a4a/papers/math-brain-2026-10-24}.
\bibitem{BrainFramework} Y.~Shmalo, \emph{Brain Framework}, public software repository, 2026,
\url{https://github.com/yspennstate/brain-framework-public/tree/46b935d26d5a2893f7790b60466387dd8da47ef5}.
\end{thebibliography}

\end{document}
